\subsection{The line resolvent in the RC-gap class}\label{sec:B3T05}
Consider residual boards on the ordinary board with $k\ge\delta_0 n$. In addition to the hypotheses of Section~\ref{sec:B3T01}, suppose that, with constants independent of $k$,
\[
 p_e\le C/k,\quad h_D\ge\eta>0,\quad
 \lambda_2(\mathscr L_p)\ge\gamma>0.
\]
Here $\mathscr L_p$ is the RC Laplacian of the connected row--column graph with weights $p$. The gap assumption concerns the weights and goes beyond a condition on the support alone. We do not prove here that every relevant residual board satisfies these hypotheses. The number of lines satisfies $m\le6n-2\le6k/\delta_0$.

For large $k$, we have $p_e<1/4$. Set $P=\operatorname{diag}p$ and
\[
 \tau_e=\frac{p_e(1-p_e)}{1-4p_e},\quad T=\operatorname{diag}\tau,
 \quad S=\operatorname{diag}(0_{RC},h_D)+A^TTA,
\]
where $A$ is the full square--line incidence matrix.
The kernel of $S$ is the line spanned by $V=(\mathbf1_R,-\mathbf1_C,0_D)^T$. The matrix $\Pi=VV^T/(2k)$ is the orthogonal projector onto this line, and $G=S^+$ is the pseudoinverse of $S$.

\paragraph{Gap of $S$.}
For line variables $(y,z)$ with $y\perp V$, since $\tau\ge p$,
\[
 (y,z)^TS(y,z)\ge
 \|P^{1/2}(A_{RC}y+A_Dz)\|_2^2+\eta\|z\|_2^2.
\]
The RC gap gives $\|P^{1/2}A_{RC}y\|^2\ge\gamma\|y\|^2$.
The two diagonal incidences per square and the diagonal masses, each at most one, give $\|P^{1/2}A_Dz\|^2\le2\|z\|^2$.
Applying $\|a+b\|^2\ge\theta\|a\|^2-\theta\|b\|^2/(1-\theta)$ with $\theta=\eta/(\eta+4)$ yields
\[
 S|_{V^\perp}\succeq c_* I,\quad
 c_*=\min\{\eta\gamma/(\eta+4),\eta/2\}>0.
\]
In particular, $\|G\|_2\le c_*^{-1}$.

\paragraph{Entrywise bounds.}
Since $\tau-p=3p^2/(1-4p)$ and each line has mass at most one, $S_{\ell\ell}=1+O(k^{-1})$ for all lines.
Two distinct lines intersect in at most one square, so $S_{\ell j}=O(k^{-1})$ if $\ell\ne j$.
Each line contains at most $k$ squares and has at most $3k$ neighbors with a nonzero entry. Thus, for $W=I-S$, we have $\|W_{\ell,*}\|_2=O(k^{-1/2})$.
The exact identity
\[
 G=I+W-2\Pi+WGW
\]
can be checked on every eigenspace of $S$, including its kernel. It requires neither $\|W\|<1$ nor a series expansion.
Since $|\Pi_{ij}|\le1/(2k)$,
\[
 |(WGW)_{ij}|\le c_*^{-1}\|W_{i,*}\|_2\|W_{j,*}\|_2,
 \qquad \boxed{G_{ij}=\mathbf1_{i=j}+O(k^{-1}).}
\]
There are $O_{\delta_0}(k)$ lines, so $\|G\|_{\infty\to\infty}=O_{\delta_0}(1)$. All constants depend only on $\delta_0,C,\eta,\gamma$.

\paragraph{Comparison with the gauge-fixed matrix.}
The subscript “gf” indicates gauge fixing as in Section~\ref{sec:B3T01}. In matrices and vectors indexed by lines, we delete the coordinate of the column whose constraint was omitted. The gauge-fixed inverse does not agree with $G$ in every RC entry, but the associated forms agree on vectors orthogonal to $V$. In particular, the incidence row $A_e$ satisfies $A_eV=0$ and
\[
 A_{e,\rm gf}S_{\rm gf}^{-1}A_{e,\rm gf}^T=A_eGA_e^T=4+O(k^{-1}).
\]
For a diagonal line $d$, we have $e_d\perp V$ and $(S_{\rm gf}^{-1})_{dd}=G_{dd}=1+O(k^{-1})$.
